\begin{figure}[t]
\centering
\resizebox{\columnwidth}{!}{% Microcoded engine (paper version of report/sekiller/zoo_motor.tex)
\begin{tikzpicture}[x=1cm,y=1cm, font=\sffamily\footnotesize,
  box/.style={draw=black!45, rounded corners=2pt, fill=white, align=center, minimum height=9mm, inner sep=3pt},
  acc/.style={box, draw=blue!60!black, fill=blue!7, line width=0.9pt},
  mem/.style={box, fill=black!4},
  ext/.style={box, draw=orange!70!black, fill=orange!8},
  ps/.style={box, draw=green!45!black, fill=green!7},
  flow/.style={-{Stealth[length=2mm]}, draw=black!55, line width=0.6pt}]
  \node[ext, minimum width=26mm] (prog) at (0.8,0) {\textbf{Program}\\[-1pt]\scriptsize 256-bit microcode};
  \node[acc, minimum width=30mm] (ctl) at (4.9,0) {\textbf{Controller}\\[-1pt]\scriptsize fetch / decode / dispatch};
  \node[ext, minimum width=44mm] (rom) at (10.4,0) {\textbf{Model memories (BRAM)}\\[-1pt]\scriptsize weights, $m_0$/$sh$/bias, tables, LUTs};
  \node[ps, minimum width=34mm] (ps7) at (10.4,1.55) {\textbf{Zynq PS (ARM)}\\[-1pt]\scriptsize AXI4-Lite: load, run, read};
  \draw[flow] (prog) -- (ctl);
  \draw[flow, draw=green!45!black, line width=0.9pt] (ps7) -- (rom);
  \draw[flow, draw=green!45!black, line width=0.9pt] (ps7.west) -| (prog.north);
  \foreach \n/\x/\t/\s in {emb/-0.2/EMB/{embeddings}, norm/2.2/NORM/{LN, RMS, isqrt},
                           attn/7.3/ATTN/{GQA, exp LUT}, rope/9.6/ROPE/{Q1.14 tables}, ew/12.0/Elementwise/{MUL, GRU, BoW}}
    {\node[box, minimum width=21mm, minimum height=12mm] (\n) at (\x,-2.3) {\textbf{\t}\\[-1pt]\scriptsize \s};}
  \node[acc, minimum width=23mm, minimum height=12mm] (gemv) at (4.75,-2.3) {\textbf{GEMV}\\[-1pt]\scriptsize $P$ MAC lanes,\\[-3pt]\scriptsize shadow bank};
  \draw[draw=black!55] (ctl.south) -- (4.9,-1.15);
  \draw[draw=black!55] (-0.2,-1.15) -- (12.0,-1.15);
  \foreach \n in {emb,norm,gemv,attn,rope,ew} \draw[flow] (\n.north |- 0,-1.15) -- (\n.north);
  \draw[flow, draw=blue!60!black, line width=1pt] (rom.south) ++(-1.6,0) |- (5.6,-0.75) -- (5.6,-1.6);
  \draw[draw=black!55] (-0.2,-3.55) -- (12.0,-3.55);
  \foreach \n in {emb,norm,gemv,attn,rope,ew} \draw[flow, {Stealth[length=2mm]}-{Stealth[length=2mm]}] (\n.south) -- (\n.south |- 0,-3.55);
  \node[mem, minimum width=50mm] (vrf) at (2.0,-4.6) {\textbf{Vector registers}\\[-1pt]\scriptsize 11 INT8 buffers, 3 read ports};
  \node[mem, minimum width=24mm] (rl) at (6.4,-4.6) {\textbf{R, L}\\[-1pt]\scriptsize INT16 residual, logits};
  \node[mem, minimum width=22mm] (kv) at (9.2,-4.6) {\textbf{KV cache}\\[-1pt]\scriptsize 8-byte head words};
  \node[mem, minimum width=20mm] (st) at (11.7,-4.6) {\textbf{State}\\[-1pt]\scriptsize sums, history};
  \foreach \m in {vrf,rl,kv,st} \draw[flow, {Stealth[length=2mm]}-{Stealth[length=2mm]}] (\m.north |- 0,-3.55) -- (\m.north);
\end{tikzpicture}
}
\caption{The microcoded engine. A controller fetches 256-bit instructions and dispatches them to one of seven units;
all model-specific content lives in memories that the Zynq PS loads at run time.}
\label{fig:engine}
\end{figure}

\textbf{Instruction set.} Table~\ref{tab:isa} lists the fifteen operations. A 256-bit instruction carries an opcode,
destination and two source buffers, a mode, four 16-bit sizes, two memory base addresses and two constants
($m_0$/$sh$, the normalization constants $K$ and $E$, or the softmax constant $C$ and shift $F$). The compiler emits the
program of a model (2 to 74 instructions) and its memory images; the controller executes the program once per token.

\begin{table}[t]
\centering\footnotesize
\caption{Engine operations.}
\label{tab:isa}
\setlength{\tabcolsep}{3pt}
\resizebox{\columnwidth}{!}{%
\begin{tabular}{@{}ll@{}}
\toprule
Operation & Semantics \\ \midrule
\code{EMB16}/\code{EMB8}/\code{EMBWIN} & embedding (+ position) / window of $k$ embeddings \\
\code{TABLOGIT}, \code{HRESET} & bigram logits; $h\leftarrow h_0$ at $t=0$ \\
\code{NORM} & LN/RMS over segments, INT16 or INT8 in/out \\
\code{GEMV} & $W[x;x']+b$; INT8 out, residual add, table, logits \\
\code{ROPE}, \code{KVSTORE}, \code{ATTN} & rotation; cache write; GQA softmax-average \\
\code{MUL}, \code{GRUUPD}, \code{COPY}, \code{REQ}, \code{BOW} & element-wise operations \\ \bottomrule
\end{tabular}}
\end{table}

\textbf{Units.} The GEMV unit is output-stationary with $P$ MAC lanes ($P=8$ unless stated): it streams one
$P$-byte weight word per cycle, accumulates $P$ outputs in INT32, and hands them to a shadow bank so that the serial
requantizer of block $b$ overlaps the accumulation of block $b+1$. A layer of $n_\mathrm{in}\times n_\mathrm{out}$ takes
$n_\mathrm{in}\lceil n_\mathrm{out}/P\rceil$ cycles; the input may be the concatenation of two buffers ($[x_t;h_{t-1}]$).
The NORM unit makes one pass, normalizes $Q$, takes a 28-cycle integer square root and one division. The attention unit has
8 lanes (the head size): a KV-cache word holds one head at one position, so a score and a weighted-value accumulation take
one cycle per cached position; a head costs $2(t+1)$ cycles plus one reciprocal. Divisions use a restoring divider that
starts at bit $Q_B$ when the quotient is known to be $<2^{Q_B}$ (41 instead of 64 cycles for normalization, 18 instead
of 34 for the softmax reciprocal). Element-wise units read up to three INT8 vectors per cycle.

\textbf{Cycle model.} Because every unit is a fixed state machine, the cycles of the token at position $t$ follow
from the microcode in closed form:
\begin{equation}
\begin{aligned}
C(t) ={}& 3 + 2N_\mathrm{instr} + \textstyle\sum_\mathrm{NORM} S\,(2n + 73) + C_\mathrm{other}\\
&+ \textstyle\sum_\mathrm{GEMV}\big[n_\mathrm{in} + (B-1)\max(n_\mathrm{in},P{+}2) + P + 5\big]\\
&+ \textstyle\sum_\mathrm{ATTN} H\,\big(2(t+1) + 30\big),
\end{aligned}
\label{eq:cycles}
\end{equation}
with $B=\lceil n_\mathrm{out}/P\rceil$ blocks per layer, $S$ normalization segments of length $n$, $H$ heads, and small
linear terms for the remaining operations ($C_\mathrm{other}$: embeddings $N{+}2$, RoPE $2+3d_h/2$ per head, element-wise
$n$ or $n{+}3$). Eq.~\eqref{eq:cycles} reproduces the simulated cycles of every model at every position exactly
(\zooPerfPairs{} (run, position) pairs, maximum deviation \zooPerfMaxErr{} cycles), so the hardware cost of an
architecture can be read off its program without simulation or synthesis.

\textbf{Storage.} Eleven INT8 vector buffers (P-byte words, byte-enable writes, three asynchronous read ports), the INT16
residual and logit memories, the KV cache, the running sums of the bag-of-words model and the token history are on chip.
Weights, requantization parameters, tables and activation tables are block RAMs.

\textbf{One bitstream.} With \code{LOADABLE=1} the model memories are RAMs written through a load port; an AXI4-Lite wrapper
exposes the load port, per-step control, a logit window, a hardware argmax and a cycle counter to the ARM cores of the
Zynq. The \emph{universal} configuration is sized for the maximum over all models ($V=512$, $T=256$, five layers); a
Python driver loads any model in milliseconds to seconds and runs it without re-synthesis. With \code{SPECIALIZE=1} the
units a model does not use are removed, which we use to attribute resources to architectural features.
